%% transmon-benchmark.13 -- revision of v12, 2026-10-15, against the Opus and
%% Fable reviews of v12, the operator's rulings on framing, Section 8 and the
%% derivative, and the follow-up runs of
%% benchmarks/transmon_devicelayout_controls/ (PRIMA in our environment, a
%% Palace-only Newton, a second refinement family on the step-2 design). See
%% changelog.md, provenance.md and validation.md. v1-v12 and their critic
%% siblings are unchanged.

\documentclass{anvil-paper}
\usepackage{xurl}
\usepackage{array}

%% --- URL line-breaking: allow breaks at hyphens and add stretch so the
%% bibliography stays inside the measure (render-gate overfull check).
\expandafter\def\expandafter\UrlBreaks\expandafter{\UrlBreaks\do\-}
\Urlmuskip=0mu plus 1mu
%% Float placement: let text share pages with floats.
\renewcommand{\topfraction}{0.85}
\renewcommand{\bottomfraction}{0.6}
\renewcommand{\textfraction}{0.12}
\renewcommand{\floatpagefraction}{0.75}
\setcounter{topnumber}{3}
\setcounter{totalnumber}{4}
%% Compact bibliography.
\renewcommand{\bibfont}{\footnotesize}
\setlength{\bibsep}{1pt plus 1pt}

%% --- Draft-phase helper macros -------------------------------------------
%% Figure placeholder: uses the rendered figure if pub-figures has produced
%% it, otherwise a framed placeholder naming the source script.
\newcommand{\anvilfig}[3][0.98]{%
  \IfFileExists{figures/#2}%
    {\includegraphics[width=#1\linewidth]{figures/#2}}%
    {\fbox{\parbox{#1\linewidth}{\centering\vspace{5em}%
      \texttt{figures/#2}\\[0.5em]%
      {\small placeholder: render via \texttt{figures/src/#3}}%
      \vspace{5em}}}}}

\title{Uniform mesh refinement underestimates the frequencies of the
DeviceLayout.jl transmon: an independent finite-element cross-check and
re-optimization}
\author{Robb Walters\\
  {\normalsize Independent Researcher}}
% NOTE: Crutcher Dunnavant to be offered coauthorship at review (operator, 2026-10-13)
\date{October 2026}

\begin{document}
\maketitle

\begin{abstract}
The SingleTransmon example of DeviceLayout.jl publishes the qubit and
resonator frequencies of an initial and an optimized transmon design,
computed with Palace. We reproduce them with Palace to within $0.7\%$ and
cross-check Palace with geode-fem, an independently implemented
finite-element solver; on shared meshes the two agree to $10^{-4}\%$ on the
qubit and to $0.03\%$ on the resonator. Under refinement the qubit
frequencies keep rising. On the initial design the qubit of an adaptive
Palace run reaches $4.4297$\,GHz, $7.0\%$ above the published $4.14$\,GHz,
with no extrapolation. Extrapolated at order 1 in the mesh size, the
fastest rate the sheet-edge singularity allows and so an estimate that
leans low, the two qubits lie $7$--$9\%$
and the resonators $2.7\%$ above their published values, and a certified
bracket on the island capacitance, read through a circuit model, is
consistent with that shift. Re-optimized for the $3.0$/$4.0$\,GHz targets at
the extrapolated frequencies, the design needs a capacitor about
$19$--$20\%$ longer than published (about $+17\%$ to $+39\%$ within the
extrapolation band). Newton's method in DeviceLayout's variables meets its
stopping criterion in 3 Palace runs, with geode-fem's shape derivative or a
scaling law in the Jacobian, where DeviceLayout's PRIMA optimizer takes 10
in the same environment.
\end{abstract}

\section{Introduction}
\label{sec:intro}

Superconducting-qubit layouts are designed by iterating field solves
against frequency targets, a loop that QDesignOptimizer
\citep{eriksson2025automated} steers with analytic models. DeviceLayout.jl \citep{devicelayout} builds the layout and
its mesh, and Palace \citep{palace}, an MFEM-based solver \citep{mfem},
computes the eigenmodes. DeviceLayout's SingleTransmon example
\citep{devicelayout-blog,devicelayout-example} runs the loop end to end. It
publishes an initial transmon-plus-resonator design at $4.14$ and
$5.591$\,GHz and a design optimized for targets of $3.0$ and $4.0$\,GHz, at
$3.005$ and $3.982$\,GHz, found with a derivative-free optimizer in nine
Palace runs.

Those frequencies come from one discretization, and the example's authors
say what that implies. They run Palace at second order ``to reduce
sensitivity to the discretization'' and note that this ``does not mean that
the eigenfrequencies found by Palace will have converged to their true
values'' \citep{devicelayout-example}. This paper measures
how far from converged the published frequencies are. Our
main result is that the published frequencies are discretization-limited.
Palace's adaptive refinement of the initial design reaches a qubit $7.0\%$ above the
published value without extrapolation, and its uniform ladders extrapolate
to qubits $7$--$9\%$ above both published values. geode-fem reproduces every
level of those ladders it can reach to $10^{-4}\%$ on the qubit, and a
certified bound on the island capacitance, read through a circuit model, is
consistent with that shift and caps the initial qubit under that model.

The second code is GEODE-FEM (``geode-fem''), a Rust finite-element solver
on the Burn tensor framework \citep{burn}. It solves the same
$H(\mathrm{curl})$ eigenproblem as Palace with a separate implementation,
and it computes design derivatives, among them the shape derivative of an
eigenfrequency. On a mesh of DeviceLayout's initial design, given as the
same input file to both codes, geode-fem reproduces Palace's six eigenmodes above $4.5$\,GHz to
$\leq 0.033\%$ and the qubit mode to $5\times10^{-5}\,\%$
(Section~\ref{sec:agreement}). This agreement establishes that the two codes
assemble and solve the same discrete problem. It also gives an independent
same-mesh cross-check of the Palace/DeviceLayout.jl transmon stack,
complementing the earlier checks of Palace by SQDMetal
\citep{sommers2025open} and \citet{ye2025electromagnetic}.

Our contributions are:
\begin{itemize}
  \item \textbf{A refinement study of a published design example.}
    Resolution ladders at first and second order put the extrapolated
    qubits $7.3\%$ and $8.7\%$ above the published values and the
    resonators $2.7\%$ above; a second refinement family and an adaptive
    run check the qubit estimate (Section~\ref{sec:crosscheck}).
  \item \textbf{Certified capacitance brackets with the sapphire tensor.}
    An equilibrated-flux lower bound and a conforming upper bound bracket
    the island capacitance of the unmodified device in $[79.56,
    79.76]$\,fF; under a two-node circuit model they cap the initial qubit
    at $4.477$\,GHz (Section~\ref{sec:static}).
  \item \textbf{A re-optimized design, with controls.} At the extrapolated
    frequencies the $3.0$/$4.0$\,GHz design needs $c \approx
    1451$--$1466\,\mu$m against the published $1217.678\,\mu$m. Newton's
    method meets DeviceLayout's criterion in 3 Palace runs with the shape
    derivative or a scaling law in its Jacobian, the latter also with no
    geode-fem input at all; Broyden's method takes 5 and PRIMA 10 in the
    same environment (Section~\ref{sec:reopt}).
  \item \textbf{Same-input-mesh agreement with Palace} on six eigenmodes
    and the qubit mode, with the $17.49$\,GHz mode that both codes share
    traced to the distributed port model (Section~\ref{sec:agreement}).
\end{itemize}

Section~\ref{sec:discussion} sorts the results by the evidence behind them
and explains why the shape derivative gives no advantage on this benchmark.

\section{Related Work}
\label{sec:related}

\paragraph{Cross-checks of Palace.}
SQDMetal \citep{sommers2025open} cross-checks a Palace workflow against
COMSOL and Ansys with mesh-convergence studies, and
\citet{ye2025electromagnetic} compare a Palace workflow with cryogenic
measurement at the ${\sim}0.3\%$ level. Both use models built
independently for each tool. Our comparison instead holds the input mesh
and the junction model fixed and changes the solver implementation, and it
adds a refinement study of a published design: code verification
\citep{oberkampf2010verification}, not validation against measurement. As of October 2026, neither Palace nor DeviceLayout.jl has a
methods publication; we cite their repositories and documentation.

\paragraph{Design sensitivities and qubit-design optimization.}
Adjoint sensitivity analysis of frequency-domain electromagnetic solvers is
an established field \citep{nikolova2004adjoint}, including shape adjoints
for Maxwell cavity eigenvalues \citep{akcelik2005adjoint}; Palace exposes no
derivatives.
In qubit design, the SQcircuit line \citep{rajabzadeh2023analysis} and
SuperGrad \citep{wang2025supergrad} differentiate at the circuit-Hamiltonian
level, SQuADDS \citep{shanto2024squadds} replaces simulation by lookup in a
validated design database, and \citet{eun2023shape} shape-optimize transmon
pads with FEM and a derivative-free global optimizer.
DeviceLayout's example optimizes with PRIMA \citep{zhang2023prima}, a
reference implementation of Powell's model-based derivative-free methods,
which build an interpolation model of the objective from function values;
with bounds only it runs BOBYQA \citep{powell2009bobyqa}.

\paragraph{Certified two-sided capacitance bounds.}
Bracketing a capacitance between complementary energy principles is
classical \citep{specogna2011complementary}. Our lower bound uses the
hypercircle idea of \citet{prager1947approximations} with Raviart--Thomas
fluxes \citep{raviart1977mixed} equilibrated patch by patch
\citep{braess2007equilibrated,ern2015polynomial}. The method is not new; we
apply it to an island capacitance with a floating conductor on a
production qubit layout, and we have found no certified capacitance bound in the
qubit-design literature.

\section{Solvers, Device and Models}
\label{sec:setup}

\subsection{geode-fem}
\label{sec:architecture}

geode-fem assembles first-order N\'ed\'elec \citep{nedelec1980} and scalar
P1/P2 element operators as batched $[n_{\mathrm{elem}},k,k]$ tensor
operations on Burn. Its eigensolver is shift-invert Lanczos \citep{arpack}
on a real symmetric pencil, factorized with faer \citep{kazdadi2026faer}. The two codes share the
formulation, the lumped junction port, the materials and the input mesh,
and geode-fem follows a specification of Palace's operators
(Section~\ref{sec:built}). Assembly, the eigensolver and the linear algebra
are separate implementations, so agreement tests the implementation and
not the shared modelling choices.

The eigenfrequency derivative used in Section~\ref{sec:reopt} needs no
adjoint solve, because the pencil is symmetric: for a simple eigenpair
$(\lambda, x)$ with $x^{\mathsf T}Mx = 1$,
\begin{equation}
  \frac{d\lambda}{d\theta} \;=\;
  x^{\mathsf T}\Bigl(\frac{\partial K}{\partial\theta}
  - \lambda\,\frac{\partial M}{\partial\theta}\Bigr)x ,
  \label{eq:hf-eig}
\end{equation}
the classical eigenvalue derivative of a symmetric generalized eigenproblem
\citep{foxkapoor1968rates}. Here $K$ and $M$ include the junction terms of
Eq.~\eqref{eq:port}; the junction does not move under the capacitor-length
motion of Section~\ref{sec:reopt}, so those terms do not contribute.
geode-fem evaluates $\partial K/\partial\theta$ and $\partial
M/\partial\theta$ along a node velocity $\partial X/\partial\theta$ with a
hand-written forward-mode dual-number twin of its tensor-permittivity
N\'ed\'elec kernel, whose real part is tested against the assembled kernel
to $10^{-10}$ relative. Burn's autodiff tape is not used for any derivative
in this paper.

\subsection{Device, mesh and materials}
\label{sec:geometry}

The device is DeviceLayout.jl v1.15.0's SingleTransmon example
\citep{devicelayout}: a transmon with a meander readout resonator and a
feedline on sapphire, in a vacuum box (Figure~\ref{fig:geometry}). Its two
design parameters are the transmon's \texttt{capacitor\_length} $c$ and the
resonator's \texttt{total\_length} $L$; the published initial design has
$c = 620\,\mu$m and $L = 5000\,\mu$m. The sapphire permittivity is
$\mathrm{diag}(9.3, 9.3, 11.5)$ with the optic axis normal to the chip;
the configuration rotates the material axes about that normal, which
leaves a uniaxial tensor unchanged, so the lab tensor is diagonal. Where a
scalar model is used it is the trace average $\bar\varepsilon = 10.03$.

\begin{figure}[htbp]
  \centering
  \anvilfig{fig1-geometry.pdf}{fig1\_geometry.py}
  \caption{Benchmark geometry, drawn from the SHA-256-pinned MSH~4.1
  input mesh of the DeviceLayout.jl v1.15.0 \texttt{SingleTransmon} device.
  (a) Top view of the three metal conductors of Section~\ref{sec:static}:
  the feedline (teal), the qubit island (blue), and the ground plane with
  the meander resonator (gray); squares mark the two readout ports. (b) The
  $1\times1\,\mu$m junction \texttt{lumped\_element} sheet with its
  four-triangle triangulation. (c) Material cross-section, schematic and not
  to scale. Both solvers read this file (22{,}684 nodes; 133{,}314
  tetrahedra).}
  \label{fig:geometry}
\end{figure}

Two kinds of mesh appear. The \emph{fixture} is an earlier generator's mesh
of the initial geometry (22{,}684 nodes, 133{,}314 tetrahedra), pinned by
SHA-256;
Section~\ref{sec:agreement} and the capacitance brackets of
Section~\ref{sec:static} use it and its refinements, which leave its
polygonal surfaces unchanged. For the cross-check and the re-optimization,
DeviceLayout meshes each geometry with the two parameters exposed, at first
order for the ladders and at its own second order for the runs that
reproduce its discretization. Its second-order mesh has the same
tetrahedra with curved faces; at the same element order, curving changes
the qubit by at most $0.0001\%$ and the resonator by about $0.05\%$. On the
initial geometry the fixture and DeviceLayout's mesh give the same Palace
qubit at $h$ ($3.4527$ against $3.4520$\,GHz, $0.02\%$). PEC is applied on
the metal sheets and the box; the readout ports are left open in every
eigenproblem except the DeviceLayout-configuration runs of
Section~\ref{sec:crosscheck}.

\subsection{Junction and circuit models}
\label{sec:junction}

The Josephson junction is a linear lumped element, $L_J = 14.860$\,nH in
parallel with $C_J = 5.5$\,fF, on the four-triangle sheet $\Gamma$ of
Figure~\ref{fig:geometry}b, as in DeviceLayout's Palace configuration. With
port length $\ell$ and width $w$ the uniform lumped-port construction adds
\begin{equation}
  K_{\mathrm{port}} = \frac{\ell}{w\,L_J}\, S_\Gamma
  \qquad\text{and}\qquad
  M_{\mathrm{port}} = \frac{C_J\,\ell}{w}\, S_\Gamma
  \label{eq:port}
\end{equation}
to the pencil $(K + K_{\mathrm{port}})x = \omega^2(M + M_{\mathrm{port}})x$,
where $S_\Gamma$ is the isotropic surface matrix on $\Gamma$. Palace's
\texttt{LumpedPort} adds the same terms with the same scalings. geode-fem's
rank-1 port-voltage shunt replaces $S_\Gamma$ by $b\,b^{\mathsf T}$, with
$b_i = (1/w)\int_\Gamma N_i\cdot\hat e\,dS$ the port-voltage functional; it
keeps the physical modes to four decimals on the fixture, removes the
artifact mode of Section~\ref{sec:agreement}, and is the junction model of
every geode-fem eigensolve in Sections~\ref{sec:crosscheck}
and~\ref{sec:reopt}.

The electrostatic model splits the metal into three conductors: ground
with the resonator, the floating feedline, and the island. Its island
capacitance with the feedline floating and neutral,
$C_\Sigma = C_{ii} - C_{if}^2/C_{ff}$, excludes $C_J$. A two-node circuit
connects it to the eigenproblem: a bare qubit
$1/(2\pi\sqrt{(L_J + L_s)(C_\Sigma + C_J)})$, with $L_s$ the series
inductance of the junction's current path, coupled with coefficient $k$
to the resonator. On the fixture $L_s$ and $k$ were measured on three and
two discretizations; extrapolated, they give bands of $0.113$--$0.125$\,nH
and $0.0227$--$0.0310$, which are estimates.

\section{Same-Input-Mesh Agreement With Palace}
\label{sec:agreement}

Both codes read the fixture with the same materials and junction model at
first order (Table~\ref{tab:agreement}). Palace's eigensolver searches above
its target (read from its source at the commit of
Section~\ref{sec:repro}), so the six modes above $4.5$\,GHz come from a run
with \texttt{Target}~$=4.5$\,GHz and the qubit from one with
\texttt{Target}~$=2$\,GHz. All six agree to $\leq 0.033\%$ (worst case
$0.032\%$), and the qubit mode to $5\times10^{-5}\,\%$, with port
participation $0.9942$ in geode-fem and port EPR $0.99424$ in Palace.

\begin{table}[htbp]
  \centering
  \caption{Eigenmode agreement from the identical input mesh (22{,}684
  nodes, 133{,}314 tetrahedra), first order in both solvers, junction as
  the lumped $L_J \parallel C_J$ shunt. geode-fem solves $156{,}863$
  N\'ed\'elec unknowns ($133{,}108$ interior); Palace refines the mesh
  internally around the metal sheets and solves $183{,}891$. geode-fem
  frequencies are given to the archived three decimals (six for the qubit),
  and $\Delta$ is the archived relative difference. Source:
  \texttt{benchmarks/transmon\_eigen/results.toml}.}
  \label{tab:agreement}
  \footnotesize
  \begin{tabular}{@{}llccc@{}}
    \toprule
    Mode & & geode-fem (GHz) & Palace (GHz) & $\Delta$ (\%) \\
    \midrule
    qubit & transmon qubit & 3.452751 & 3.452749306805 & 0.00005 \\
    resonator & readout resonator & 5.153 & 5.151335830348 & 0.032 \\
    mode 2 & box/cavity & 15.465 & 15.46052107794 & 0.029 \\
    port sheet & port-sheet artifact & 17.490 & 17.49010903536 & 0.001 \\
    mode 4 & box/cavity & 18.693 & 18.69165792915 & 0.007 \\
    mode 5 & box/cavity & 20.703 & 20.69755679425 & 0.026 \\
    mode 6 & box/cavity & 26.088 & 26.08089940472 & 0.027 \\
    \bottomrule
  \end{tabular}
\end{table}

The $17.49$\,GHz mode is a shared artifact of the isotropic sheet term:
the junction patch has one free node whose hat-function gradient carries no
port voltage but still stores energy in $S_\Gamma$, with Rayleigh quotient
$1/(L_JC_J)$, or $f_{LC} = 17.60$\,GHz, lowered to $17.49$\,GHz by the
adjacent volume mass; Palace gives it a port EPR of $2.49\times10^{-8}$.
This mechanism is our inference, supported by a synthetic test; the rank-1
shunt removes the mode, and its $0.001\%$ agreement tests only the port
operator.

The fixture is coarse. Palace's own refinements of it move the qubit from
$3.4527$\,GHz to $3.9528$\,GHz at $h/2$ ($+14.5\%$) and to $4.1721$\,GHz
at second order on $h$ ($+20.8\%$), and the resonator by $+6.1\%$ at
$h/2$. The agreement above is code verification on one discretization;
Section~\ref{sec:crosscheck} measures the discretization error.

\section{The DeviceLayout Cross-Check}
\label{sec:crosscheck}

\paragraph{The published designs.}
The DeviceLayout.jl blog \citep{devicelayout-blog} reports the initial
design ($c = 620\,\mu$m, $L = 5000\,\mu$m) at $[4.14, 5.591]$\,GHz and
the final design ($c = 1217.678\,\mu$m, $L = 7243.955\,\mu$m) at
$[3.005, 3.982]$\,GHz, reached by PRIMA in nine Palace runs ``within 1\% of
our target frequencies''. The example source \citep{devicelayout-example}
sets the targets to $[3.0, 4.0]$\,GHz, works in the variables $x_1 =
\sqrt{620\,\mu\mathrm{m}/c}$ and $x_2 = 5000\,\mu\mathrm{m}/L$, and stops
when $\mathrm{err} = \frac12\sum_m (f_m - t_m)^2/t_m^2 < 10^{-4}$. Its
frequencies are computed with Palace on DeviceLayout's mesh at mesh order 2
and solver order 2, with a first-order absorbing far field and two
$50\,\Omega$ readout ports.

\paragraph{Reproduction at DeviceLayout's discretization.}
Palace on DeviceLayout's second-order mesh with the example's
configuration file gives $[4.1683, 5.6132]$\,GHz for the initial design
and $[3.0164, 3.9956]$\,GHz for the final one, $0.68\%$ and $0.40\%$
above the published initial pair and $0.38\%$ and $0.34\%$ above the
final pair. The same cells with the PEC-box model used everywhere else
differ from the configuration-file runs by at most $0.003\%$.

\paragraph{The ladder and its rate.}
For each geometry, Palace runs at order 1 on DeviceLayout's first-order
mesh $h$ and its uniform refinements $h/2$ and $h/4$ (gmsh's eight-way
split, up to $12.6$ million unknowns), and at order 2 on $h$ and $h/2$;
geode-fem runs at order 1 on $h$ and $h/2$, on the same meshes. Every
computed series rises under refinement (Figure~\ref{fig:crosscheck}); this
is an observation, not a guarantee. At the edge of a zero-thickness metal
sheet the field behaves as $E \sim r^{-1/2}$ \citep{meixner1972behavior},
so near the edge the eigenfunctions lie in $H^{s}$ only for $s < 1/2$, and
on uniform meshes the eigenvalue error decays as $h^{2s}$ at any element
order \citep{boffi2010finite}, slower than $O(h)$. The order-1
ladders show pre-asymptotic orders of $0.61$ (qubit) and $0.70$--$0.73$
(resonator).

\paragraph{Extrapolation rule.}
The \emph{estimate} is the mean of two extrapolations at order 1 in the
frequency, $2f_2 - f_1$: one from the order-1 pair ($h/2$, $h/4$) and one
from the order-2 pair ($h$, $h/2$). Order 1 is the fastest rate the
theory allows, so the rule makes the smallest correction consistent with
it, and its estimates lean low; for the order-2 pair the order is taken
from theory, because two levels give no observed order. The \emph{band}
spans six values: those two; the same pairs extrapolated at order 1 in the
eigenvalue variable $\lambda\propto f^2$, $\sqrt{2f_2^2 - f_1^2}$, which
set the low side; and the same pairs at the observed order, which set the
high side. The band is an extrapolation range, not a bound. For the
initial qubit, the circuit-model cap of Section~\ref{sec:static} trims the
band at $4.477$\,GHz. The rule was fixed after the static route had been
computed: extrapolation at the observed order alone, the rule used before,
puts the initial qubit at $4.63$\,GHz, above the circuit-model cap, and the
observed-order values are kept only as the band's high side. The agreement
of the two routes is therefore a consistency check, not a blind
prediction.

\begin{figure}[htbp]
  \centering
  \anvilfig{fig2-crosscheck.pdf}{fig2\_crosscheck.py}
  \caption{Resolution ladders of the two published DeviceLayout.jl
  SingleTransmon designs: Palace at orders 1 and 2 on DeviceLayout's mesh
  $h$ and its uniform refinements, and geode-fem at order 1 on the same
  meshes (crosses); the open diamond is Palace at DeviceLayout's own
  discretization, that of the published values (dashed).
  The ``$h\to0$'' column is categorical: the extrapolated estimate with its
  band and, for the initial geometry, the estimate of the second refinement
  family (open triangle). In (a), the static route of
  Section~\ref{sec:static} and the circuit-model cap. Sources:
  \texttt{benchmarks/transmon\_devicelayout\_crosscheck/results.toml} and
  \texttt{benchmarks/transmon\_devicelayout\_controls/results.toml}.}
  \label{fig:crosscheck}
\end{figure}

\begin{table}[htbp]
  \centering
  \caption{Published frequencies, Palace at DeviceLayout's discretization,
  and the extrapolated estimates with their bands, in GHz, from the first
  refinement family. Percentages are relative to the published value.
  Source: \texttt{benchmarks/transmon\_devicelayout\_crosscheck/results.toml}.}
  \label{tab:crosscheck}
  \footnotesize
  \begin{tabular}{@{}lcccc@{}}
    \toprule
    & Published & At DeviceLayout's discretization & Extrapolated estimate & Band \\
    \midrule
    Initial qubit & 4.14 & 4.1683 ($+0.68\%$) & 4.444 ($+7.3\%$) & $[4.415, 4.677]^{\ast}$ \\
    Initial resonator & 5.591 & 5.6132 ($+0.40\%$) & 5.740 ($+2.7\%$) & $[5.723, 5.817]$ \\
    Final qubit & 3.005 & 3.0164 ($+0.38\%$) & 3.265 ($+8.7\%$) & $[3.236, 3.471]$ \\
    Final resonator & 3.982 & 3.9956 ($+0.34\%$) & 4.088 ($+2.7\%$) & $[4.077, 4.139]$ \\
    \bottomrule
  \end{tabular}

  \smallskip
  {\footnotesize $^{\ast}$Under the circuit model of
  Section~\ref{sec:static}, the cap of $4.477$\,GHz trims the band to
  $[4.415, 4.477]$. The second refinement family gives $4.472$ $[4.458,
  4.510]$.}
\end{table}

\paragraph{Result.}
The initial qubit extrapolates to $4.444$\,GHz and the final qubit to
$3.265$\,GHz, $7.3\%$ and $8.7\%$ above the published values; both
resonators extrapolate to $2.7\%$ above theirs
(Table~\ref{tab:crosscheck}). Every published value lies below its band.
Measured against the estimates, DeviceLayout's discretization sits $6.2\%$
and $7.6\%$ low on the two qubits and $2.2\%$ and $2.3\%$ low on the
resonators. With no extrapolation, the finest uniform levels already
lie $4.1\%$ and $4.7\%$ above the published qubit values. The estimates rest
on Palace's order-1 $h/4$ and order-2 levels, which geode-fem did not run.
On the levels it did run, geode-fem and Palace agree on the qubit to
$10^{-4}\,\%$ on both geometries. On the resonator they differ by $0.03\%$
at $h$, where Palace adds $4{,}946$ (initial) and $5{,}944$ (final)
elements by local bisection, and by $0.0000\%$ at $h/2$, where both codes
solve the same mesh.

\paragraph{A second refinement family.}
Three checks test the initial-geometry estimate, and all three concern the
qubit (Figure~\ref{fig:initialq}). The first is a second refinement family:
Palace's own uniform refinement of the same $h$ mesh, which first bisects
the elements at the metal sheet. Its $h/2$ qubit is $3.9469$\,GHz against
gmsh's $3.8771$\,GHz, and on the fixture the two refinements give $3.9528$
and $3.8939$\,GHz, so the refinement scheme accounts for $1.5$--$1.8\%$ of
the $h/2$ level and the choice of $h$ mesh for only $0.15$--$0.43\%$. The
rule gives $4.472$ $[4.458, 4.510]$\,GHz for the qubit, at an observed order
of $0.91$ against the first family's $0.61$, and $5.776$ $[5.768,
5.800]$\,GHz for the resonator. The two estimates lie $0.64\%$ apart and the
second band lies inside the first, so we quote the initial qubit as
$4.444$--$4.472$\,GHz, $7.3$--$8.0\%$ above the published value.

\paragraph{Adaptive refinement.}
A further check comes from Palace's adaptive mesh refinement, run from the
same $h$ mesh at order 2, with the error indicator summed over the qubit and resonator modes,
D\"orfler marking at $0.7$ and nonconformal refinement. In five solves from
$0.95$ to $11.0$ million unknowns the qubit rises from $4.1683$ to
$4.4297$\,GHz, $0.32\%$ below the estimate and inside the band. The third
solve, at $2.2$ million unknowns, already exceeds the finest uniform level
($4.3118$\,GHz at $7.0$ million). The run has not reached a plateau (the
last step moved the qubit by $+0.55\%$; $3.1$\,h, $93$\,GiB peak on 32
ranks). A three-point fit $f_\infty - CN^{-\alpha}$
through the last three solves gives $4.4552$\,GHz ($\alpha = 0.80$), but it
is exactly determined and ill-conditioned: $+0.1\%$ on the finest level
moves it by $24$\,MHz. The adaptive resonator rises more slowly under this
marking, and its finest level, $5.6986$\,GHz, is still below its band, so
the adaptive run checks the qubit estimate only.

\paragraph{A common-limit fit.}
The last check fits the
order-1 ($h/2$, $h/4$) and order-2 ($h$, $h/2$) series of each ladder with
one common limit and order. This fit is exactly determined; it lies
$0.5$--$1.3\%$ above the rule's estimate on every first-family series, as
the rule's leaning low predicts, and $+0.1\%$ on the order-2 $h/2$ level
alone moves the initial-qubit fit by $20$\,MHz. Section~\ref{sec:reopt}
repeats the second family on the re-optimized design.

\begin{figure}[htbp]
  \centering
  \anvilfig{fig3-initial-qubit.pdf}{fig3\_initial\_qubit.py}
  \caption{Checks of the extrapolated estimate on the initial geometry.
  Left: every Palace level against the number of N\'ed\'elec unknowns, for
  the two refinement families of DeviceLayout's mesh $h$ (filled: gmsh's
  eight-way split, the ladder of Figure~\ref{fig:crosscheck}; open:
  Palace's refinement) and Palace adaptive refinement from $h$ at order 2
  (diamonds). Right, on the same frequency axis: the estimates of the two
  families with their bands and, for the qubit only, the adaptive run's
  ill-conditioned three-point fit; the resonator fit ($\alpha = 0.25$) is
  not reliable and not shown. In (a), the circuit-model cap. Sources:
  \texttt{benchmarks/transmon\_devicelayout\_crosscheck/results.toml} and
  \texttt{benchmarks/transmon\_devicelayout\_controls/results.toml}.}
  \label{fig:initialq}
\end{figure}

\section{A Static Route: Certified Island Capacitance}
\label{sec:static}

\paragraph{The certificate.}
For the reduced electrostatic problem (island at potential 1, ground and
resonator at 0, feedline floating with zero net charge), the potential
$\phi$ satisfies $0 \le \phi \le 1$ by the maximum principle, and
$C_\Sigma = \int \varepsilon\,\nabla\phi\cdot\nabla\phi$ with $\varepsilon
= \varepsilon_0\varepsilon_r$. Every conforming finite-element energy is an
upper bound on $C_\Sigma$, because $C_\Sigma$ is the minimum of the energy
over admissible potentials. For the lower bound, take any piecewise
polynomial field $D$ on the mesh cut along the zero-thickness sheets.
Integrating by parts and using $0\le\phi\le1$ on every term gives
\begin{equation}
  \int D\cdot\nabla\phi \;\ge\; N(D) \;=\; Q_{\mathrm{island}}(D)
  - \lvert Q_{\mathrm{feed}}(D)\rvert
  - \lVert \operatorname{div} D\rVert_{L^1}
  - \lVert [D\cdot n]\rVert_{L^1(\mathcal F)},
  \label{eq:cert-n}
\end{equation}
where $Q$ are the fluxes of $D$ into the two conductors and
$[D\cdot n]$ is the normal jump on the interior faces $\mathcal F$ off
the sheets. Cauchy--Schwarz, $\bigl(\int D\cdot\nabla\phi\bigr)^2 \le
\int D\cdot\varepsilon^{-1}D \int\varepsilon\nabla\phi\cdot\nabla\phi$,
then gives, whenever $N(D) > 0$,
\begin{equation}
  C_\Sigma \;\ge\; \frac{N(D)^2}{\int D\cdot\varepsilon^{-1}D}.
  \label{eq:cert}
\end{equation}
The bound holds for any $D$, so solver tolerance can only weaken it. We
take $D$ from the P2 potential, equilibrated patch by patch into
Raviart--Thomas fluxes of order one to three, with an independent
Crouzeix--Raviart flux as a check. The construction is validated against
an exact parallel plate and an analytic coaxial bracket, and mutation tests
reject a grounded-feedline or jump-free certificate. \emph{Certified} in this paper means rigorous
for the discrete polygonal geometry of the mesh, with lower bounds
floored and upper bounds ceiled. The certificate is evaluated in plain
f64, not interval arithmetic; the worst-case summation error is about
$10^{-10}$ relative, far below the rounding. Refinement does not move the
polygonal surfaces, so all meshes of one geometry bound the same exact
capacitance, and a family bracket is the largest lower bound and the
smallest upper bound.

\paragraph{The unmodified device.}
Uniform refinement of the fixture with P1 and P2 elements gives the upper
bounds of Figure~\ref{fig:convergence}: with the scalar model, from
$137.706752$\,fF at P1 on $h$ to $83.015285$\,fF at P2 on $h/2$, and with
the tensor from $136.684731$ to $85.008425$\,fF. On adaptively refined meshes the
certificate closes the gap: the island capacitance of the unmodified
initial design lies in $[77.40, 77.61]$\,fF with the scalar
$\bar\varepsilon$ and in $[79.56, 79.76]$\,fF with the sapphire tensor
(Figure~\ref{fig:convergence}). The certified tensor-minus-scalar
difference, computed from the unrounded bounds, is $1.95$--$2.35$\,fF, or
$2.5$--$3.1\%$ of the scalar value. The fixture itself overstates the
tensor $C_\Sigma$ by at least $56.92$\,fF, which is at least $71.3\%$ of
the exact value, and the scalar one by at least $77.4\%$.

\begin{figure}[htbp]
  \centering
  \anvilfig{fig4-convergence.pdf}{fig4\_convergence.py}
  \caption{Mesh convergence of the island capacitance $C_\Sigma$ of the
  unmodified device under uniform refinement of the fixture, with the scalar trace-averaged
  permittivity (blue) and the sapphire tensor (orange). Every marker is a
  conforming upper bound on the exact capacitance of the meshed polygonal
  geometry. The short bars are the certified two-sided brackets from
  adapted meshes, each drawn at the tetrahedron count of the mesh that
  sets its lower bound ($3{,}289{,}543$ and $3{,}388{,}197$) and magnified
  in the inset on a fF scale; the inset title gives the certified
  tensor-minus-scalar gap. Sources: \texttt{c\_sigma\_convergence.toml},
  \texttt{c\_sigma\_bracket.toml} and
  \texttt{c\_sigma\_bracket\_tensor.toml}.}
  \label{fig:convergence}
\end{figure}

\paragraph{From capacitance to the qubit frequency.}
Reading the certified tensor bracket through the two-node circuit of
Section~\ref{sec:junction}, over every corner of the $C_\Sigma$, $L_s$ and
$k$ bands, with the resonator at its extrapolated estimate, gives the
initial dressed qubit in $[4.449, 4.458]$\,GHz. That range lies inside the
eigenmode band of Section~\ref{sec:crosscheck} and inside the range of the
two family estimates, and the first family's estimate is $0.22\%$ below
its midpoint. Because the extrapolation rule was chosen after this range
was known, the $0.22\%$ is a consistency check, not an independent
prediction. The range is an estimate, because $L_s$ and $k$ are.

\paragraph{A model-dependent cap.}
Under the circuit model, a series inductance and a coupling to a higher
mode can only lower the qubit, so the certified lower bound $C_\Sigma \ge
79.56$\,fF caps it at $1/(2\pi\sqrt{L_J(C_\Sigma + C_J)}) = 4.477$\,GHz.
The cap is rigorous only given that model (one island node with $L_J
\parallel C_J$, $L_s \ge 0$, coupling only to modes above the qubit, a
quasi-static $C_\Sigma$) and rests on the fixture, not on the ladder meshes,
which agree with it to $0.02\%$ at $h$, so we list it as model-dependent,
not certified. It trims the band's observed-order high side ($4.677$\,GHz)
and does not depend on the extrapolation rule. Every computed level and both family
estimates lie below it. The first family's common-limit fit, $4.4802$\,GHz,
lies $3.6$\,MHz ($0.08\%$) above it, an excess equal to a $0.019\%$ change
of one order-2 level and so within that fit's conditioning.

\section{Re-Optimization and Its Controls}
\label{sec:reopt}

The published final design misses its targets at the extrapolated
frequencies. We repeat DeviceLayout's optimization against its own
discretization, to compare optimizers in one Palace environment, and
against the extrapolated frequencies, to find the design that meets the
targets there.

\paragraph{The capacitor\_length derivative.}
DeviceLayout keeps the junction fixed and builds the island upward from
it, so a change of $c$ translates every node above the island rigidly and
stretches the island band. geode-fem turns this motion into a node velocity
$\partial X/\partial c$, checks that every metal outline node in a
stretched band moves along its own edge ($3{,}146$ nodes on the initial
mesh), and evaluates Eq.~\eqref{eq:hf-eig} with the sapphire tensor for both modes.
The contraction takes $0.44$\,s. On the initial mesh it matches a central
difference of two re-solves of the moved mesh ($\pm0.5\,\mu$m) to every
printed digit: $-2.443227\times10^{-3}$\,GHz/$\mu$m for the qubit and
$-3.964\times10^{-6}$\,GHz/$\mu$m for the resonator.

\paragraph{The total\_length column.}
DeviceLayout re-lays the meander when $L$ changes, so no fixed-topology
node motion represents it. A transmission-line model holds $f_r(L+\delta)$
constant and takes $\partial f_q/\partial L = 0$, with $\delta =
216\,\mu$m from one finite difference of two geode-fem order-1 solves
($L = 5000$ and $5250\,\mu$m). $\delta$ is weakly determined: remesh noise
of about $0.15\%$ in either $f_r$ moves it by about $\pm175\,\mu$m. The
model enters only the Jacobian, and every iterate is evaluated on the true
objective, so a poorer Jacobian costs steps, not accuracy.

\paragraph{Newton on the high-fidelity objective.}
DeviceLayout re-meshes every design, so the first-order objective is not
smooth in the parameters: from $c = 620$ to $650\,\mu$m the remeshed finite
difference of the qubit frequency is $-5.23\times10^{-3}$\,GHz/$\mu$m,
against a fixed-topology derivative of $-2.28\times10^{-3}$. The iteration
therefore runs Newton directly on the high-fidelity objective $F$, in
DeviceLayout's variables $(x_1, x_2)$, in which the frequencies are close
to linear and nearly decoupled: a probe in either variable moves its own
mode by $4.5\%$ and the other by at most $0.05\%$. The $c$ column of the
Jacobian is geode-fem's first-order derivative scaled by the fidelity
ratio $S = F/f_h$ at the same design. Each step costs one geode-fem
first-order solve with its gradient ($19$--$25$\,s) and one high-fidelity
evaluation; the one-time $\delta$ calibration is not counted.

\begin{table}[htbp]
  \centering
  \caption{Palace runs to DeviceLayout's criterion ($\mathrm{err} <
  10^{-4}$) at its discretization, counting the initial design as run~1,
  all in our Palace environment (build, mesher and platform) except the
  published PRIMA run. The Newton variants use the transmission-line $L$
  column with the stated $\delta$; Broyden's count includes its two probes.
  Sources: the \texttt{optimize} and \texttt{controls} artifacts under
  \texttt{benchmarks/transmon\_devicelayout\_*/}.}
  \label{tab:controls}
  \footnotesize
  \begin{tabular}{@{}lllccc@{}}
    \toprule
    Method & $c$ column & $\delta$ from & Runs & First-step err & Design at criterion ($\mu$m) \\
    \midrule
    PRIMA, published & none & none & 9 & unpublished & (1217.678, 7243.955) \\
    PRIMA, our environment & none & none & 10 & n/a & (1209.125, 7207.551) \\
    Newton & shape derivative & geode-fem & 3 & $7.9\times10^{-4}$ & (1225.346, 7223.250) \\
    Newton & scaling law & geode-fem & 3 & $1.7\times10^{-4}$ & (1232.664, 7221.270) \\
    Newton & scaling law & Palace & 3 & $2.1\times10^{-4}$ & (1229.026, 7225.712) \\
    Broyden & two probes & none & 5 & $6.2\times10^{-4}$ & (1224.847, 7232.447) \\
    \bottomrule
  \end{tabular}
\end{table}

\paragraph{Against DeviceLayout's discretization.}
Table~\ref{tab:controls} and Figure~\ref{fig:reopt}(b) compare the
optimizers on one objective. Newton
with the shape-derivative column meets the criterion at its second step,
after 3 Palace runs, and its three geode-fem solves add about $70$\,s,
against about $4$\,min per Palace run on 32 ranks; a third step reaches
$\mathrm{err} = 1.4\times10^{-7}$ at $(1236.360, 7224.900)\,\mu$m. Two
gradient-free controls test where the 3 runs come from. The first replaces
the $c$ column by the scaling law $\partial F_q/\partial x_1 = F_q/x_1$,
which follows from $f_q \propto 1/\sqrt{C_\Sigma} \propto x_1$. With the
geode-fem $\delta$ it also takes 3 runs, and its first step has an err
$4.7$ times lower ($1.7\times10^{-4}$ against $7.9\times10^{-4}$), about
$2.2$ times closer in frequency. With $\delta = 563.879\,\mu$m taken from a
Palace probe instead, so that no geode-fem input remains, it again takes 3.
The second control, Broyden's method \citep{broyden1965class} from two
forward finite-difference probes, takes 5. DeviceLayout's own PRIMA call,
re-run verbatim in our environment, takes 10: runs 1--5 are its initial
interpolation set of $2n+1 = 5$ points \citep{powell2009bobyqa}, and run 9 misses the criterion by $15\%$ ($\mathrm{err} =
1.15\times10^{-4}$). In this environment the published final design already
meets the criterion ($\mathrm{err} = 1.5\times10^{-5}$), so the environment
moves PRIMA's trajectory, not the attainability of the target; the
published count of 9 is from DeviceLayout's environment. In one
environment, then, the comparison is 3 Newton runs against 10.

\paragraph{Why the shape-derivative step lands farther.}
At the initial design the scaled shape derivative gives
$\partial F_q/\partial x_1 = 3.66$\,GHz, within about $1.5\%$ of the $3.71$\,GHz
of Broyden's Palace probe, so the derivative measures the local slope
well. Both local slopes are $11$--$12\%$ shallower than the scaling law's
$F_q/x_1 = 4.168$\,GHz, because the frequency is not exactly proportional
to $x_1$ near the initial design, and the first step is long ($x_1$ from
$1$ to about $0.7$): the chord slope from the initial design to the
shape-derivative design at the criterion is $4.0$\,GHz, close to the
scaling law and steeper than either tangent. A local tangent therefore
overshoots: the
shape-derivative step lands at $c = 1338.338\,\mu$m and Broyden's at
$1320.522\,\mu$m, against about $1230\,\mu$m at the criterion, while the
scaling law, which follows $f_q \propto x_1$ over the whole step, lands at
$1196.933\,\mu$m.

\paragraph{Against the extrapolated frequencies.}
Starting from the published final design, each Newton step uses the
extrapolated estimate and the geode-fem derivative at the current design,
and each new design is evaluated by a full ladder with the rule of
Section~\ref{sec:crosscheck} (Table~\ref{tab:newton-conv}).
The first overshoots the qubit by $-1.09\%$; the second lands at $+0.16\%$
and $-0.04\%$ of the targets, with both targets inside the bands. On both
ladders geode-fem and Palace agree on the qubit to $2\times10^{-4}\,\%$.
One scaling-law step from the same start, evaluated by the Palace half of
the ladder, lands $3.6$ and $8.2\,\mu$m short of step~2 and at $+0.13\%$
and $+0.12\%$ of the targets: one gradient-free step does what the two
shape-derivative steps did. Each ladder took $1.4$ and $1.6$\,h of Palace
wall time, dominated by the order-2 $h/2$ cell.

\begin{table}[htbp]
  \centering
  \caption{Steps against the extrapolated frequencies, in GHz. Each design
  is evaluated by a full resolution ladder (the scaling-law step and the
  second family by Palace alone) and the rule of
  Section~\ref{sec:crosscheck}; values are estimates with their bands.
  Family b refines step 2's $h$ mesh with Palace's uniform refinement
  instead of gmsh's. Sources: the \texttt{optimize}, \texttt{controls} and
  (published final) \texttt{crosscheck} artifacts under
  \texttt{benchmarks/transmon\_devicelayout\_*/}.}
  \label{tab:newton-conv}
  \footnotesize
  \begin{tabular}{@{}lcccc@{}}
    \toprule
    Design & $c$ ($\mu$m) & $L$ ($\mu$m) & $f_q$ [band] & $f_r$ [band] \\
    \midrule
    Published final & 1217.678 & 7243.955 & 3.265 $[3.236, 3.471]$ & 4.088 $[4.077, 4.139]$ \\
    Step 1 & 1481.657 & 7407.945 & 2.967 $[2.938, 3.188]$ & 4.004 $[3.993, 4.055]$ \\
    Step 2 & 1446.078 & 7416.352 & 3.005 $[2.979, 3.221]$ & 3.998 $[3.987, 4.067]$ \\
    Step 2, family b & 1446.078 & 7416.352 & 3.020 $[3.003, 3.101]$ & 4.022 $[4.016, 4.043]$ \\
    Scaling-law step & 1442.479 & 7408.181 & 3.004 $[2.977, 3.235]$ & 4.005 $[3.994, 4.055]$ \\
    \bottomrule
  \end{tabular}
\end{table}

\paragraph{The re-optimized design.}
Step 2 is the design whose extrapolated frequencies sit on the targets:
$c = 1446.078\,\mu$m, $18.8\%$ longer than published, with a resonator
$2.4\%$ longer ($+172\,\mu$m). The second refinement family at step 2 gives
$3.0203$ $[3.0033, 3.1014]$\,GHz for the qubit and $4.0225$\,GHz for the
resonator, $0.52\%$ and $0.60\%$ above the first family, the same spread as
on the initial geometry; both targets lie just below its band, so on family
b step 2 is slightly short. Mapped to $3.0$\,GHz by the scaling law $f_q
\propto c^{-1/2}$, the two families' qubit estimates call for $c = 1451$ and
$1466\,\mu$m, $+19.1\%$ and $+20.4\%$ over the published length. Mapped
linearly through $dF/dc \approx
S\,\partial f_h/\partial c = -9.05\times10^{-4}$\,GHz/$\mu$m, the slope at
the step-1 design where step~2 started, and centred on the $3.0$\,GHz
target, the first family's qubit band corresponds to $c$ between about
$1423$ and $1690\,\mu$m, about $+17\%$ to $+39\%$; under the scaling law the
same band gives $+17.1\%$ to $+36.9\%$, and it contains family b's band
($+19.0\%$ to $+26.9\%$). Since the estimates lean low, the common-limit
fit, which puts step~2 at $3.039$ and $4.019$\,GHz, would call for a longer
capacitor still.

\begin{figure}[htbp]
  \centering
  \anvilfig{fig5-reoptimization.pdf}{fig5\_reoptimization.py}
  \caption{The re-optimization and its controls. (a) The design plane near
  the optimum: the published final design, PRIMA's end point in our
  environment, the Newton iterates against DeviceLayout's discretization
  (numbered: shape derivative; the first legs start off the plot), and the
  steps against the extrapolated frequencies, with the step-2 qubit band
  mapped linearly to capacitor\_length around the $3.0$\,GHz target (an
  estimate). (b) DeviceLayout's err against the number of Palace
  evaluations at its discretization, in our environment; PRIMA's published
  end point is plotted at its 9 runs. This counts Palace evaluations, not
  total cost: each shape-derivative Newton point also costs one geode-fem
  solve. (c, d) Extrapolated estimates and bands against the targets, with
  step~2 in both refinement families. Sources: the \texttt{optimize},
  \texttt{controls} and \texttt{crosscheck} artifacts under
  \texttt{benchmarks/transmon\_devicelayout\_*/}.}
  \label{fig:reopt}
\end{figure}

\section{How geode-fem Was Designed and Built}
\label{sec:built}

geode-fem implements the whiteroom L1--L4 operator specification of
Crutcher Dunnavant \citep{dunnavant2026whiteroom}. He also gave design
input on the solver's operator layer and contributed code: the migration to
Rust 2024, the Metal (Apple GPU) backend, the present layout of the core
crate and the examples, and the geode-util crate. The rest of the code, the
experiments and the text of this paper were produced by frontier AI models
working as role-separated agents: Anthropic's Claude (Opus, Fable and
Sonnet models, through Claude Code) and OpenAI's GPT Sol and GPT Astra
models (via Codex). The agents were coordinated by Loom \citep{loom}, an
issue, worktree and pull-request workflow, and by a shared Squad room.
Nearly every change was reviewed by a separate agent session. The reviews
of this paper recomputed its numbers from the archived artifacts,
reproduced runs, and mutation-tested the scripts that check the text and
the figures against those artifacts. The human author directed the
research, made every scoping and claim decision, and takes responsibility
for the content.

\section{Reproducibility}
\label{sec:repro}

Every result in this paper is read from an archived, scripted
artifact.\footnote{The geode-fem source is public at
\url{https://github.com/rjwalters/geode-fem}. Archival DOI to be added at
deposit.}
% TODO(operator): Zenodo DOI
Each table and figure caption names its artifact; the circuit terms $L_s$
and $k$ come from \path{benchmarks/transmon_loop_inductance/} and
\path{benchmarks/transmon_coupling/}, the Palace cells are under
\path{reference/fixtures/transmon_palace/}, and the DeviceLayout meshes are
generated by \path{reference/julia/generate_transmon_param.jl}. The
generated artifacts are pinned by regression tests. The versions are geode-fem commit
\texttt{3174015} for the same-mesh agreement, with the commit of each later
result recorded in its artifact, Palace \texttt{fba6a5b}, DeviceLayout.jl
v1.15.0, PRIMA.jl 0.2.4 and Burn 0.21. The DeviceLayout ladders, the
adaptive run and the controls ran on AWS Graviton4 hosts
(\texttt{r8g.8xlarge}, 32 vCPU, and \texttt{r8g.4xlarge}, 16 vCPU); the
capacitance brackets ran on a 28-core, 96\,GiB workstation. geode-fem's
direct solves at the ladders' $h/2$ peaked at $94$ and $136$\,GB, which is
why only Palace ran the larger cells. The Palace runs read an MSH~2.2
conversion of each mesh that keeps every node and physical group. The figure scripts in
\path{figures/src/} read the artifacts directly.

\section{Discussion}
\label{sec:discussion}

\paragraph{What is established, and how.}
Table~\ref{tab:tiers} sorts the results by the kind of evidence behind
them. The brackets are the only bounds, and the cap follows from a bracket
only under a circuit model. What stands without any extrapolation is the adaptive run's $4.4297$\,GHz, $7.0\%$ above the
published initial qubit, and the finest uniform levels, $4.1\%$ and $4.7\%$
above the two published qubits; under the monotone rise every series shows,
these are lower estimates.

\begin{table}[htbp]
  \centering
  \caption{Evidence tiers. Certified: a guaranteed bound for the meshed
  polygonal geometry. Model-dependent: a rigorous consequence of a certified
  bound under a stated circuit model. Measured: a solve on a stated
  discretization. Estimated: an extrapolation with a stated rule.
  Consistency: an agreement that is not independent evidence.}
  \label{tab:tiers}
  \footnotesize
  \begin{tabular}{@{}>{\raggedright\arraybackslash}p{0.17\linewidth}>{\raggedright\arraybackslash}p{0.77\linewidth}@{}}
    \toprule
    Tier & Results \\
    \midrule
    Certified & $C_\Sigma$ of the unmodified device: $[77.40, 77.61]$\,fF (scalar), $[79.56, 79.76]$\,fF (tensor); tensor-minus-scalar gap $2.5$--$3.1\%$ \\
    Model-dependent & The initial-qubit cap of $4.477$\,GHz: two-node circuit, quasi-static $C_\Sigma$, bracket certified on the fixture's polygonal geometry, not on the ladder meshes \\
    Measured & Same-mesh agreement (Table~\ref{tab:agreement}); Palace at DeviceLayout's discretization within $0.7\%$ of the published pairs; every ladder and adaptive level (adaptive initial qubit $+7.0\%$, finest uniform levels $+4.1\%$ and $+4.7\%$); the gradient check; Palace runs to DeviceLayout's criterion in one environment: 3 for each Newton variant, 5 for Broyden's method, 10 for PRIMA \\
    Estimated & Extrapolated frequencies and bands (Table~\ref{tab:crosscheck}), the initial qubit $4.444$--$4.472$\,GHz over two refinement families; the adaptive fit and the common-limit fits; the static-route qubit range; $L_s$ and $k$; the re-optimized capacitor, $+19$ to $+20\%$ over two families and about $+17$ to $+39\%$ within the band \\
    Consistency & The two routes agree to $0.22\%$ on the initial qubit; the extrapolation rule was fixed after the static route was known, so this is not a blind prediction \\
    Not claimed & Agreement with measurement; a verified hit of $3.0$/$4.0$\,GHz in the $h \to 0$ limit; a global optimum; an advantage of the shape derivative on this problem; eigenmode-participation (EPR) derivatives; a speed advantage over Palace \\
    \bottomrule
  \end{tabular}
\end{table}

\paragraph{Why the shape derivative gives no advantage here.}
The derivative is a central tool of the re-optimization, but three
properties of this benchmark remove its advantage. First, the problem has
two nearly decoupled parameters, close to linear in DeviceLayout's
variables, so Newton's method converges with any reasonable Jacobian: the
shape derivative, a two-probe secant, or the scaling law $f_q \propto x_1$,
which costs nothing. Second, a finite-difference gradient costs only $N+1 =
3$ solves. Third, the coarse-mesh, fixed-topology derivative is the slope
of a different discrete problem from the one the iteration evaluates:
DeviceLayout remeshes every design and Palace evaluates it at second
order, so a remeshed finite difference of that objective is the more
faithful slope; here the scaled derivative lies about $1.5\%$ below one. That the
derivative is exact for its own discrete problem and cheap ($0.44$\,s)
does not change the number of Palace runs.

Derivatives pay off in other regimes. With $N$ parameters, a discrete
adjoint gives all $N$ derivatives of one objective for about one extra
solve, where finite differences need $N$ more; coupled targets with no
per-parameter scaling law need a computed Jacobian; and freeform shapes
have many parameters by construction. A companion study
\citep{walters2026conformal} applies geode-fem's discrete shape adjoint to
a conformal patch antenna with 38 shape parameters; its comparison with
gradient-free search is still being run, so it does not yet show whether
the derivative helps there.

\paragraph{Threats and limitations.}
Agreement between two codes is evidence against implementation error, not
against shared modelling choices such as the port sheet behind the
$17.49$\,GHz mode. The extrapolations assume the rate the sheet-edge singularity
allows, their high sides are wide, and the adaptive run had not reached a
plateau. The circuit route uses a bracket certified on a different mesh of
the same initial geometry, and $L_s$ and $k$ were extrapolated from two or
three discretizations. The certificate is evaluated in f64 and validated by
tests, not formally verified. The estimates are for DeviceLayout's geometry
under a lossless, linear model ($L_J \parallel C_J$, open readout ports in
the ladders, zero-thickness metal), not for a fabricated device, and only
one device is studied. The eigenmode branch is
not yet differentiated with respect to participation ratios, which needs
eigenvector derivatives \citep{nelson1976simplified}.

\paragraph{What a practitioner can do.}
The slow rate comes from the sheet edges, and the standard remedy for an
edge singularity is to refine toward it, with a mesh graded toward the
edges or with $hp$ refinement \citep{demkowicz2006computing}, which recover
faster convergence than uniform refinement. Palace's adaptive run is a
measured instance: at $2.2$ million unknowns it already exceeds the finest
uniform level at $7.0$ million.

\section{Conclusion}
\label{sec:conclusion}

Palace reproduces the published frequencies of the DeviceLayout.jl
SingleTransmon example at the example's own discretization, and that
discretization sits $6$--$8\%$ low on the qubits. An adaptive Palace run
puts the initial qubit $7.0\%$ above its published value with no
extrapolation, and the extrapolated qubit-frequency estimates are $7$--$9\%$
above both published values; the zero-thickness sheet edge limits uniform
refinement to at best $O(h)$, so these estimates lean low. geode-fem
reproduces Palace's qubit on every level both codes ran, and a certified
bracket on the island capacitance, read through a circuit model, is
consistent with that shift and caps the initial qubit under that model. At
the extrapolated frequencies the $3.0$/$4.0$\,GHz design needs a capacitor
about $19$--$20\%$ longer than published, and Newton's method in
DeviceLayout's variables reaches the example's criterion in 3 Palace runs
against PRIMA's 10, with or without the shape derivative. For design
examples that are
meant to be copied, the discretization is part of the result and belongs
in the report beside it.

\paragraph{Acknowledgments.}
We thank Crutcher Dunnavant for the whiteroom specification, for his
design input to the solver's operator layer and for his contributions to
the codebase, and the Palace and DeviceLayout.jl developers for publishing
the SingleTransmon example and its frequencies. The AI-orchestrated
development of geode-fem and of this paper is described in
Section~\ref{sec:built}.
% TODO(operator): approve acknowledgment wording

\bibliographystyle{plainnat}
\bibliography{refs}

\end{document}
